% Replacement for Appendix A only. Keep the existing \appendix command.
% Stop before the section labelled app:complexity. Do not append to old A.
\section{\textcolor{blue}{Reversible construction of the point-arithmetic circuit}}
\label{app:oracle}

\begin{revision}
The construction uses affine-coordinate arithmetic and controlled point
translations within the framework of \cite{ProosZalka2003,Roetteler2017}.
The elliptic-curve group law is standard
\cite{Silverman2009,Washington2008}. The explicit identity flag, total
extension to invalid encodings, branch organisation, and arithmetic
assembly described below specify the reference implementation used here.
They are not presented as improvements over the optimized circuits of
the cited studies, and their resource estimates are analysed separately.

Let $\ell=\lceil\log_2 p\rceil$ and use the point encoding in
Equation~\eqref{eq:point_encoding}. On valid point encodings, the scalar
evaluation has the action
\begin{equation}
\label{eq:reversible_circuit}
\begin{aligned}
U_{f_P}:\quad
&\lvert k\rangle\lvert\mathrm{enc}(R)\rangle
 \lvert0\rangle_{\mathrm{aux}}\\
&\longmapsto
 \lvert k\rangle\lvert\mathrm{enc}(R+kP)\rangle
 \lvert0\rangle_{\mathrm{aux}}.
\end{aligned}
\end{equation}
Initializing $R=\mathcal O$ gives the periodic-function evaluation in
Equation~\eqref{eq:oracle}. The auxiliary-zero requirement is part of
each block's interface.
\end{revision}

\subsection{\textcolor{blue}{Classical precomputation}}
\begin{revision}
For the scalar-register width $m$ selected in Algorithm~\ref{alg:hybrid},
compute
\[
P_0=P,\qquad P_{j+1}=2P_j,
\qquad j=0,\ldots,m-2.
\]
Thus $P_j=2^jP$, using at most $m-1$ classical point doublings. These
constants configure the controlled translations. A constant
$P_j=\mathcal O$ contributes no update and is omitted. Neither the order
of $P$ nor a table of all $2^m$ scalar outputs is required. The values
$A$, $-A$, and $2A$ needed by a finite constant-point block are also
classical constants, rather than quantum inputs to be enumerated.
\end{revision}

\subsection{\textcolor{blue}{Reversible modular arithmetic}}
\begin{revision}
The reference blocks are assembled from $X$, controlled-$X$, Toffoli,
and multi-controlled-$X$ operations. Standard reversible gate
constructions provide the underlying framework
\cite{Barenco1995,NielsenChuang2010}; the particular assembly and
workspace management below describe the local implementation.

\textbf{Comparisons and flags.}
Equality and range predicates are written into flag qubits by XOR,
preserving their input registers. Equality with a constant is implemented
by temporarily complementing input bits corresponding to zeros in that
constant, applying a multi-controlled $X$ to the flag, and undoing the
complements. Range comparisons use reversible subtraction in workspace
to extract the comparison result before reversing the subtraction.
Flags are cleared by recomputing the appropriate predicate; they are
not measured or discarded during arithmetic evaluation.

\textbf{Addition.}
For a classical constant $a\in\{0,\ldots,p-1\}$, the addition block maps
\begin{equation}
\label{eq:app_r5_constant_addition}
\begin{aligned}
&\lvert x\rangle\lvert0\rangle_{\mathrm{aux}}\\
&\longmapsto
 \lvert(x+a)\bmod p\rangle
 \lvert0\rangle_{\mathrm{aux}},\quad x<p.
\end{aligned}
\end{equation}
Encodings with $x\geq p$ remain unchanged. A validity flag records
$x<p$. For valid inputs, addition is performed in an extended register
and $p$ is subtracted if the sum is at least $p$. If
$u=(x+a)\bmod p$, reduction occurred exactly when $u<a$, allowing the
reduction flag to be cleared using the output. Repeating the range
comparison clears the validity flag because modular addition preserves
the valid range and invalid inputs were fixed.

For $y=\sum_{j=0}^{\ell-1}y_j2^j$, addition of a quantum register is
assembled by letting $y_j$ control addition of $2^j\bmod p$. This maps
$(x,y)$ to $((x+y)\bmod p,y)$ for $x<p$. Subtraction is obtained by
reversing the corresponding addition circuit. Every constituent block,
including its cleanup operations, is controlled when required.

\textbf{Multiplication and squaring.}
Multiplication uses a separate modular accumulator:
\begin{equation}
\label{eq:app_r5_multiply_accumulate}
\begin{aligned}
&\lvert x\rangle\lvert y\rangle\lvert t\rangle
 \lvert0\rangle_{\mathrm{aux}}\\
&\longmapsto
 \lvert x\rangle\lvert y\rangle
 \lvert(t+xy)\bmod p\rangle
 \lvert0\rangle_{\mathrm{aux}},\quad t<p.
\end{aligned}
\end{equation}
The identity $xy=\sum_{i,j}x_i y_j2^{i+j}$ gives an implementation
through constant additions of $2^{i+j}\bmod p$, jointly controlled by
$x_i$ and $y_j$. Both operands are preserved. Invalid accumulator
encodings are fixed. Squaring uses the analogous expansion of $x^2$
and accumulates its value without overwriting $x$. A zero accumulator
therefore yields a product or square, while reversal subtracts it.

\textbf{Modular inversion.}
The inversion evaluator preserves its input and XORs its inverse into
an arbitrary target. For an $\ell$-bit input $w$, define
\[
I_p(w)=
\begin{cases}
w^{-1}\bmod p,&0<w<p,\\
0,&\text{otherwise}.
\end{cases}
\]
The clean interface is
\begin{equation}
\label{eq:app_r5_inverse_evaluator}
\begin{aligned}
&\lvert w\rangle\lvert t\rangle
 \lvert0\rangle_{\mathrm{aux}}\\
&\longmapsto
 \lvert w\rangle\lvert t\oplus I_p(w)\rangle
 \lvert0\rangle_{\mathrm{aux}}.
\end{aligned}
\end{equation}
Here $\oplus$ denotes bitwise XOR, not modular addition.

The implemented evaluator uses quotient-based Euclidean updates,
initialized with $(r_0,r_1,c_0,c_1)=(p,w,0,1)$. In each active round,
\begin{equation}
\label{eq:app_r5_euclidean_update}
\begin{aligned}
q&=\lfloor r_0/r_1\rfloor,\\
(r_0,r_1)&\longmapsto(r_1,r_0-qr_1),\\
(c_0,c_1)&\longmapsto(c_1,c_0-qc_1\bmod p).
\end{aligned}
\end{equation}
The invariant $r_i\equiv c_iw\pmod p$ is preserved. For a valid
nonzero input, termination at $(r_0,r_1)=(1,0)$ gives
$c_0=w^{-1}\bmod p$. The circuit uses $2\ell+1$ scheduled rounds,
with activation flags disabling updates after $r_1=0$. This bound
follows from the reduction of the Euclidean remainder by at least
a factor of two within every two active rounds.

The compact variant keeps one remainder/coefficient quartet and
retains quotient and activation history for reversal. The inverse is
XORed into the target under the validity and terminal-state predicates;
the complete computation is then reversed. This clears the Euclidean
state, history, and arithmetic workspace. The method is the local
reference implementation, not an adoption of the optimized inversion
circuit or resource bounds of \cite{Roetteler2017}.
\end{revision}

\subsection{\textcolor{blue}{Controlled point translation}}
\begin{revision}
For a finite constant $A=(x_A,y_A)$ and an ordinary input
$R=(x_R,y_R)$ with $x_R\neq x_A$, evaluate the affine group law
\cite{Silverman2009,Washington2008}:
\begin{equation}
\label{eq:affine_translation}
\begin{aligned}
\lambda&=(y_R-y_A)(x_R-x_A)^{-1}\pmod p,\\
x_{R+A}&=\lambda^2-x_R-x_A\pmod p,\\
y_{R+A}&=\lambda(x_R-x_{R+A})-y_R\pmod p.
\end{aligned}
\end{equation}
The implementation computes the coordinate differences, their required
inverse, and the slope, retaining shared intermediates until both output
coordinates have been XORed into their targets. The arithmetic
computation is then reversed once to clear those intermediates.

\textbf{Validity and exceptional cases.}
A finite encoding is valid only when its coordinates are canonical
residues and satisfy the curve equation. The identity is valid only
with its identity flag set and both coordinate registers zero.
Mutually exclusive branch flags select the following outputs:
\[
R+A=
\begin{cases}
A,&R=\mathcal O,\\
\mathcal O,&R=-A,\\
2A,&R=A\text{ and }A\neq-A,\\
\text{ordinary affine sum},&\text{otherwise, for valid }R.
\end{cases}
\]
When $A=-A$, the opposite-point branch takes precedence, so
$A+A=\mathcal O$. The first three outputs are constant encodings;
the ordinary branch alone invokes the nonzero-denominator formula.
Validity predicates use the curve equation, not a list of curve points.

\textbf{Total evaluator and in-place cleanup.}
Let $H_A$ map valid encodings to the encoding of their translation by
$A$, and fix invalid encodings. Translation preserves the valid set and
has inverse translation by $-A$, so $H_{-A}=H_A^{-1}$ on all encodings.
The evaluator has the clean XOR action
\begin{equation}
\label{eq:app_r5_total_point_evaluator}
\begin{aligned}
&\lvert z\rangle\lvert t\rangle
 \lvert0\rangle_{\mathrm{aux}}\\
&\longmapsto
 \lvert z\rangle\lvert t\oplus H_A(z)\rangle
 \lvert0\rangle_{\mathrm{aux}}.
\end{aligned}
\end{equation}
In particular, an invalid input is XORed unchanged into the target;
the evaluator does not leave an arbitrary target untouched.

Starting with a zero target, evaluate $H_A$, exchange the point and
target registers, and evaluate $H_{-A}$ into the old-input register.
Since $H_{-A}(H_A(z))=z$, the old encoding cancels by XOR. This is
the construction in Equation~\eqref{eq:point_update_cleanup}, with
the zero target serving as workspace rather than as a representation
of $\mathcal O$. Controlling all three stages, including exchange and
cleanup, gives $V_A$. For control zero the registers are unchanged;
for control one the point is translated and all auxiliary registers
return to zero. Invalid point encodings are fixed by the completed
in-place update.

The compact evaluator shares arithmetic intermediates and reuses a
scratch pool between branch selection and ordinary arithmetic, whose
scratch lifetimes do not overlap. Branch flags remain available until
their final uncomputation. These are local implementation choices,
not claimed improvements over published optimized arithmetic.
\end{revision}

\subsection{\textcolor{blue}{Coherent scalar evaluation}}
\begin{revision}
Initialize the accumulator in $\mathrm{enc}(\mathcal O)$ and apply
$V_{P_j}$ with control $k_j$ for $j=0,\ldots,m-1$. After the first
$j$ bits, the accumulator represents
\begin{equation}
\label{eq:partial_scalar_accumulation}
R_j=\sum_{i=0}^{j-1}k_iP_i,\qquad R_0=\mathcal O.
\end{equation}
The next update gives $R_{j+1}=R_j+k_jP_j$, and hence $R_m=kP$.
The scalar bits are preserved and the auxiliary workspace is reused
only after each update has cleared it. The basis-state action extends
by linearity to a superposition of scalars. No scalar-output lookup
table or row-dependent enumeration is used.

The point-register width is $2\ell+1$, distinct from the scalar width
$m$. Preparing the identity requires setting its flag to one while
leaving the coordinate registers zero. The all-zero point-register
string is not the chosen identity encoding. Fourier transformation,
measurement, and probabilistic recovery belong to the surrounding
procedure in Algorithm~\ref{alg:hybrid}, not to this arithmetic block.
\end{revision}

\subsection{\textcolor{blue}{Implementation correspondence and validation}}
\begin{revision}
The arithmetic contracts correspond to the functions
\texttt{modular\_add\_constant},
\texttt{modular\_add\_register},
\texttt{modular\_multiply\_accumulate}, and
\texttt{modular\_square\_accumulate}.
The compact inversion is implemented by
\texttt{modular\_inverse\_inplace}, using
\texttt{binary\_quotient} and retained Euclidean history;
the name refers to in-place internal state updates, while its external
interface remains the clean XOR evaluator above.

Point validity and branch selection correspond to
\texttt{point\_encoding\_check} and
\texttt{point\_addition\_cases}.
The shared-intermediate ordinary evaluator is
\texttt{ordinary\_point\_sum\_shared}; the total compact evaluator is
\texttt{point\_sum\_evaluator\_compact}.
The evaluate--exchange--cleanup construction and its controlled
version correspond to \texttt{point\_add\_inplace} and
\texttt{controlled\_point\_add}, selected with
\texttt{evaluator\_variant="compact"}.
The \texttt{EllipticOrderFinding} class assembles these updates with
register preparation and the inverse QFT.

Validation compares arithmetic outputs against modular and
elliptic-curve reference calculations and checks input preservation,
exceptional branches, invalid-encoding behaviour, and zero auxiliary
outputs. Since the point-arithmetic circuits consist of positive-phase
computational-basis permutations, their basis-state action also
determines their coherent action on the tested initialized subspace.
Complete order-finding validation propagates the actual arithmetic
gate sequences and computes the Fourier measurement distribution
from their resulting outputs, rather than assuming an ideal
$k\mapsto kP$ table. Such simulation validates noiseless functional
behaviour; it is not a hardware experiment or a fault-tolerant resource
demonstration. Quantitative recovery and signature results are reported
separately from these arithmetic contracts.

The gate and workspace costs of this construction, together with the
cost of repeated order-recovery executions, are analysed in
Appendix~\ref{app:complexity}.
\end{revision}
